% References Figure~\ref{fig:highlights:delve_result}, defined in the summary figure near the top of researchStatement.tex.
\textbf{Standard Siren Constraints on the Hubble Constant:}
While recent observing campaigns have not yet yielded an \EM\ counterpart to a \GW\ merger, I have retained a focus on constraining cosmology using \EM\ and \GW\ data, employing the \textit{dark siren} method to measure the local expansion rate of the universe, \Ho. 
The dark sirens method infers redshift information of \GW\ events using a catalog of potential host galaxies, which is lower precision than the bright siren method, but can be applied to all GW events \cite{2023JCAP...12..023G}.
My work began by supporting the preparation of \DES\ data, resulting in the first dark siren constraint using a deep photometric catalog to jointly infer \GW\ population parameters and cosmological parameters \citep{McMahon2026DESY6H0}. \DES\ was accepted as the fiducial result for the most recent \LVK\ analysis, and demonstrated a 22\% improvement in constraining \Ho\ compared to previous analyses \citep{LVK2026GWTC5}. I have since expanded on the \DES\ preparation, leading a dark-siren analysis using \DELVE, where initial results show that when photometric redshift biases are properly treated, \textbf{\DELVE\ outperforms \DES\ using dark sirens} (see Figure~\ref{fig:highlights:delve_result}), constraining \Ho\ to 11\% when combined with the bright siren GW170817 \citep{MacBride2026DELVEH0}. % DES+DELVE
